% !TEX program = xelatex
% 2026 CCF CCSP 竞赛出行手册
% 编译：latexmk -xelatex CCSP2026.tex   或   xelatex CCSP2026.tex（跑两遍出目录）
\documentclass[11pt,a4paper,fontset=windows]{ctexart}

\usepackage[margin=2.1cm,top=2.4cm,bottom=2.2cm]{geometry}
% 注意：不能用 xcolor 的 table 选项（会加载 colortbl）。
% 本机 MiKTeX 的 array.sty(2023) 与 colortbl.sty(2026/05 v1.0l) 版本错配，
% \rowcolor 会调用新版 array 才有的内部宏而报错。表格改用 booktabs 描线 + 彩色粗体表头。
\usepackage{xcolor}
\usepackage{booktabs}
\usepackage{tabularx}
\usepackage{longtable}
\usepackage{array}
\usepackage{multirow}
\usepackage{makecell}
\usepackage{enumitem}
\usepackage{amssymb}      % \checkmark \square
\usepackage{titlesec}
\usepackage{fancyhdr}
\usepackage{tcolorbox}
\tcbuselibrary{skins,breakable}
\usepackage[colorlinks=true,linkcolor=main,urlcolor=accent,bookmarksnumbered]{hyperref}
\usepackage{graphicx}

% ---------- 配色 ----------
\definecolor{main}{HTML}{1F4E79}      % 主色：深蓝
\definecolor{mainlite}{HTML}{2E75B6}  % 主色浅
\definecolor{accent}{HTML}{C00000}    % 强调：红
\definecolor{warnbg}{HTML}{FDF1E7}    % 警示底
\definecolor{okbg}{HTML}{EBF5EE}      % 提示底
\definecolor{infobg}{HTML}{EAF1F8}    % 信息底
\definecolor{rowalt}{HTML}{F4F8FC}    % 表格隔行
\definecolor{graytx}{HTML}{595959}

% ---------- 表格 ----------
\newcolumntype{L}[1]{>{\raggedright\arraybackslash}p{#1}}
\newcommand{\thd}[1]{\textbf{\color{main}#1}}   % 表头：主色粗体（不用 colortbl 底色）
\newcommand{\subhd}[1]{\textbf{\color{graytx}#1}}
\newcommand{\warnmark}{\textcolor{accent}{\rule[0.12em]{0.52em}{0.52em}}}  % 替代 ⚠（拉丁字体无此字形）

% ---------- 标题样式 ----------
\titleformat{\section}{\normalfont\Large\bfseries\color{main}}{\thesection}{0.7em}{}[\vspace{-0.5em}\color{main}\rule{\textwidth}{1.1pt}]
\titleformat{\subsection}{\normalfont\large\bfseries\color{mainlite}}{\thesubsection}{0.6em}{}
\titlespacing{\section}{0pt}{1.5em}{0.6em}
\titlespacing{\subsection}{0pt}{1.1em}{0.4em}

% ---------- 页眉页脚 ----------
\pagestyle{fancy}
\fancyhf{}
\fancyhead[L]{\small\color{graytx}2026 CCF CCSP 出行手册}
\fancyhead[R]{\small\color{graytx}合肥 $\rightarrow$ 成都}
\fancyfoot[C]{\small\color{graytx}\thepage}
\renewcommand{\headrulewidth}{0.4pt}

% ---------- 自定义盒子 ----------
\newtcolorbox{keybox}[1][]{
  colback=infobg, colframe=main, boxrule=0pt, leftrule=3.5pt,
  arc=1.5pt, left=8pt, right=8pt, top=6pt, bottom=6pt, #1}
\newtcolorbox{warnbox}[1][]{
  colback=warnbg, colframe=accent, boxrule=0pt, leftrule=3.5pt,
  arc=1.5pt, left=8pt, right=8pt, top=6pt, bottom=6pt, #1}
\newtcolorbox{okbox}[1][]{
  colback=okbg, colframe=green!45!black, boxrule=0pt, leftrule=3.5pt,
  arc=1.5pt, left=8pt, right=8pt, top=6pt, bottom=6pt, #1}

% 逐日行程表用的小标题
\newcommand{\daytitle}[2]{%
  \vspace{0.7em}
  {\large\bfseries\color{white}\colorbox{main}{\strut\ \ #1\ \ \normalsize #2\ \ }}%
  \vspace{0.35em}\par}

\begin{document}

% ============================================================
% 封面
% ============================================================
\thispagestyle{empty}
\vspace*{2.8cm}
\begin{center}
  {\color{mainlite}\rule{\textwidth}{2pt}}\\[0.85em]
  {\Huge\bfseries\color{main} CCSP 2026}\\[0.45em]
  {\LARGE\bfseries 出\hspace{0.35em}行\hspace{0.35em}手\hspace{0.35em}册}\\[0.85em]
  {\color{mainlite}\rule{\textwidth}{2pt}}
\end{center}

\vspace{2.4cm}
\begin{center}
\renewcommand{\arraystretch}{1.5}
\begin{tabular}{@{}r@{\hspace{2.2em}}l@{}}
\subhd{文档名称}     & 2026 CCF CCSP 竞赛出行手册\\[0.15em]
\subhd{文档版本}     & v1.0\\[0.15em]
\subhd{编\hspace{1em}制} & XCPC 集训队\\[0.15em]
\subhd{编制日期}     & 2026 年 10 月 9 日\\[0.15em]
\subhd{最近更新}     & 2026 年 10 月 9 日\\[0.15em]
\subhd{赛事名称}     & 2026 CCF 大学生计算机系统与程序设计竞赛\\[0.15em]
\subhd{主办承办}     & 中国计算机学会（CCF）主办 · 西华大学承办\\[0.15em]
\subhd{赛程日期}     & 2026 年 10 月 21 日 至 10 月 22 日\\[0.15em]
\subhd{往返日期}     & 2026 年 10 月 20 日 至 10 月 23 日\\
\end{tabular}
\end{center}

\vfill
\begin{center}
{\color{mainlite}\rule{0.45\textwidth}{0.6pt}}\\[0.6em]
{\small\color{graytx} 编制单位：合肥工业大叔队 \quad 编制日期：2026 年 10 月 9 日}
\end{center}
\newpage

% ============================================================
% 目录
% ============================================================
\thispagestyle{empty}
\tableofcontents
\newpage

% ============================================================
\section{一页速查}
% ============================================================

\begin{keybox}
\textbf{\color{main}三件最重要的事}

\vspace{0.4em}
\textbf{1. 比赛时间}：10 月 21 日 \textbf{09:00--21:00}，共 \textbf{12 小时}，中途不离场。

\vspace{0.25em}
\textbf{2. 比赛地点}：西华大学（郫都校区），成都市郫都区红光大道 9999 号。

\vspace{0.25em}
\textbf{3. \warnmark 今年新规：}
\href{https://www.ccf.org.cn/ccsp/tzgg/2026-09-14/932132.shtml}%
{\textbf{禁止携带 U 盘等任何电子存储设备}}（详见 \ref{sec:rules} 节）。
\textbf{去年允许，今年不允许！}
\end{keybox}

\vspace{0.8em}
\begin{center}
\begin{tabularx}{\textwidth}{@{}L{3.1cm}L{5.4cm}X@{}}
\toprule
\thd{项目} & \thd{内容} & \thd{说明}\\
\midrule
比赛日 & 10 月 21 日（周三）09:00--21:00 & 西华大学（郫都校区）\\
颁奖典礼 & 10 月 22 日（周四）14:00--16:00 & 天府国际会议中心\\
入住酒店 & 繁洲花园大酒店（四星） & 红光街道广场路北三段 88 号\\
酒店电话 & 028-87922176 / 15308236858 & 建议提前存手机\\
距赛场 & 约 1 公里 & 步行约 15 分钟，可打车\\
去程车次 & G3391 \quad 10:21 $\rightarrow$ 19:01 & 合肥南 $\rightarrow$ 成都南\\
返程车次 & G1888 \quad 11:51 $\rightarrow$ 21:41 & 成都东 $\rightarrow$ 合肥\\
成都气温 & 约 13--19\,$^\circ$C{}，多阴雨 & 早晚温差大，见 \ref{sec:remind} 节\\
\bottomrule
\end{tabularx}
\end{center}

\begin{okbox}
\textbf{随身必带四件套}：身份证 \quad 学生证 \quad 准考证 \quad 纸质参考资料。\\
准考证于\textbf{竞赛前 3 个工作日}开放打印，请提前打印纸质版并多备一份。
\end{okbox}

\newpage

% ============================================================
\section{竞赛信息}
% ============================================================

\subsection{基本安排}

\begin{center}
\begin{tabularx}{\textwidth}{@{}L{2.9cm}X@{}}
\toprule
\thd{项目} & \thd{内容}\\
\midrule
竞赛名称 & 2026 CCF 大学生计算机系统与程序设计竞赛（CCSP）\\
主办单位 & 中国计算机学会（CCF）\\
承办单位 & 西华大学\\
比赛日期 & 2026 年 10 月 21 日（周三）\\
比赛时间 & 09:00--21:00，共 12 小时，中午不休息、不离场\\
比赛地点 & 四川省成都市金牛区土桥金周路 999 号 西华大学（郫都校区）\\
\multirow{2}{*}{颁奖典礼} & 时间：10 月 22 日（周四）14:00--16:00\\
 & 地点：天府国际会议中心（成都市双流区蜀州路 3333 号，西博城地铁站 E 口步行 60 米）\\
赛题设置 & 全国赛共 5 题，总分 500 分，每题 100 分\\
可用语言 & C/C++、Java 或 Python\\
竞赛方式 & 机考\\
奖项 & 取全国总排名前 60\%，金银铜奖比例 1:2:3\\
附加权益 & 可免费参加 10 月 21--24 日在成都举办的 CNCC2026 大会\\
\bottomrule
\end{tabularx}
\end{center}

\begin{keybox}
\textbf{地址说明}：西华大学郫都校区在官方通知中写作「金牛区土桥金周路 999 号」，
这是历史行政区划遗留的写法；学校现用地址为\textbf{「郫都区红光大道 9999 号」}，
两者指向同一处校园。打车报地址时\textbf{建议直接用「西华大学郫都校区」}，司机更认这个。
\end{keybox}

\subsection{今年新规：物品携带}\label{sec:rules}

\begin{warnbox}
\textbf{\color{accent}\warnmark 与去年不同，请务必看完这一节}

\vspace{0.5em}
\textbf{依据（点击查看官方原文）：}\\
\href{https://www.ccf.org.cn/ccsp/tzgg/2026-09-14/932132.shtml}%
{\textbf{《关于调整 2026 CCF CCSP 竞赛选手可携带资料规定的通知》}}\\
{\small\color{graytx} CCF 官方发布 · 2026 年 9 月 14 日 · 中国计算机学会 CCSP 竞赛委员会}

\vspace{0.7em}
\textbf{\color{accent}核心调整：禁止携带任何电子存储设备（U 盘、移动硬盘、存储卡等）；
允许携带合理数量的纸质参考资料（书籍、打印材料、手写笔记等）。}

\vspace{0.7em}
\textbf{通知原文要点}
\begin{itemize}[leftmargin=1.4em,itemsep=0.28em,topsep=0.35em]
  \item \textbf{禁止}携带任何电子存储设备，包括但不限于 U 盘、移动硬盘、存储卡等。
  \item \textbf{允许}携带合理数量的纸质参考资料，包括书籍、打印材料、手写笔记等。
  \item \textbf{仍禁止}携带任何具有计算或通信功能的设备（笔记本电脑、手机、平板电脑、智能穿戴设备等）。
  \item 键盘、鼠标、有线耳机等不具备存储、计算与通信功能的外设\textbf{仍可携带}，
        入场时须经监考人员检查确认。
  \item 此前发布的竞赛章程、竞赛手册中相关条款与本通知不一致的，\textbf{以本通知为准}。
\end{itemize}
\end{warnbox}

\vspace{0.4em}
\begin{center}
\begin{tabularx}{\textwidth}{@{}L{2.85cm}L{1.55cm}X@{}}
\toprule
\thd{类别} & \thd{能否携带} & \thd{说明}\\
\midrule
电子存储设备 & \textcolor{accent}{\textbf{禁止}} & U 盘、移动硬盘、存储卡等，\textbf{任何}形式都不行\\
纸质参考资料 & \textbf{允许} & 合理数量；书籍、打印材料、手写笔记\\
计算通信设备 & \textcolor{accent}{\textbf{禁止}} & 笔记本电脑、手机、平板电脑、智能穿戴设备\\
键鼠、有线耳机 & \textbf{允许} & 须无存储 / 计算 / 通信功能，入场经检查\\
\bottomrule
\end{tabularx}
\end{center}

\newpage

% ============================================================
\section{行程总览}
% ============================================================

\begin{center}
\begin{tabularx}{\textwidth}{@{}L{2.5cm}L{2.5cm}X@{}}
\toprule
\thd{日期} & \thd{关键节点} & \thd{行程主线}\\
\midrule
10 月 20 日（周二） & 10:21 发车 / 19:01 抵达 & 合肥 $\rightarrow$ 成都南站 $\rightarrow$ 地铁 $\rightarrow$ 入住酒店\\
10 月 21 日（周三） & 09:00--21:00 & 酒店 $\rightarrow$ 西华大学，\textbf{正式比赛 12 小时} $\rightarrow$ 酒店\\
10 月 22 日（周四） & 14:00--16:00 & 酒店 $\rightarrow$ 天府国际会议中心，\textbf{颁奖典礼} $\rightarrow$ 酒店\\
10 月 23 日（周五） & 11:51 发车 / 21:41 抵达 & 酒店 $\rightarrow$ 成都东站 $\rightarrow$ 合肥\\
\bottomrule
\end{tabularx}
\end{center}

\begin{keybox}
\textbf{住宿跨度}：10 月 20 日、21 日、22 日共三晚，10 月 23 日上午退房。\\
\textbf{总时长}：10 月 20 日 10:21 出发 $\rightarrow$ 10 月 23 日 21:41 抵达合肥，全程约 3 天半。
\end{keybox}

% ============================================================
\section{逐日详细行程}
% ============================================================

以下时间均为\textbf{计划时间}。地铁、有轨电车的运行时间按正常平峰估算，未包含等车与换乘步行时间，
\textbf{当天请以高德地图 / 百度地图实时导航为准}。\\[0.3em]

% ---------- Day 1 ----------
\newpage
\daytitle{Day 1}{10 月 20 日（周二）· 合肥 $\rightarrow$ 成都}

\begin{center}
\begin{tabularx}{\textwidth}{@{}L{1.6cm}L{4.4cm}L{3.5cm}X@{}}
\toprule
\thd{时间} & \thd{事项} & \thd{交通方式} & \thd{地点 / 备注}\\
\midrule
08:45 & 集合，清点人员与行李 & — & 翡翠湖大北门\\
08:55 & 步行前往地铁站 & 步行 & 工大翡翠湖校区站（3 号线），约 \textbf{15 分钟}\\
09:10 & 乘地铁 3 号线 & 地铁 3 号线 & 前往换乘站\\
09:35 & 换乘 4 号线 & 地铁 4 号线 & \textbf{图书馆站}站内换乘\\
09:55 & 抵达合肥南站 & — & 出站后按「高铁进站口」指示走\\
10:00 & 安检、进站、候车 & — & \textbf{检票口 10A}，21 分钟后发车\\
10:21 & \textbf{G3391 发车} & 高铁 & 合肥南 $\rightarrow$ 成都南\\
\midrule
19:01 & 抵达成都南站 & 高铁 & 准点到达\\
19:05 & 出站，前往地铁站 & 步行 & 站内连通「火车南站」，跟随指示牌\\
19:18 & 乘 7 号线（内环） & 地铁 7 号线 & 往一品天下方向，乘 9 站\\
19:40 & \textbf{一品天下站换乘} & 站内换乘 & 换 2 号线，往\textbf{犀浦}方向\\
19:45 & 乘 2 号线 & 地铁 2 号线 & 乘 7 站\\
20:00 & 抵达天河路站 & — & 换乘有轨电车（需出站，留足时间）\\
20:05 & 乘有轨电车蓉 2 号线支线 & 有轨电车 & 往\textbf{仁和}方向\\
20:28 & 抵达金土站 & 有轨电车 & 酒店最近站点\\
20:35 & \textbf{抵达繁洲花园大酒店} & 步行 & 办理入住\\
20:50 & 放下行李、稍作休息 & — & 与前台确认早餐时间与退房时间\\
21:00 & 「夜宵」 & 步行 & 放松心态应赛\\
22:30 & 回房休息 & — & 次日为 12 小时比赛，\textbf{务必早睡}\\
\bottomrule
\end{tabularx}
\end{center}

\begin{warnbox}
\textbf{车上历时 8 小时 40 分}：午餐、晚餐均在车上解决，\textbf{请提前备好干粮与饮品}。
高铁停站一般只有三五分钟，来不及下车买东西；车上餐食又贵又难吃，别指望。
\end{warnbox}

\begin{keybox}
\textbf{备选方案（若地铁停运或行李过多）}\\
成都南站打车直达酒店：约 25 公里，40--50 分钟，约 70--90 元。\\
两人分摊后与地铁差距不大，\textbf{若列车晚点导致地铁赶不上末班，直接打车}。
\end{keybox}

\begin{warnbox}
\textbf{末班车提醒，务必随时看表，别到了才发现没车}\\
三者末班时间差别很大，\textbf{有轨电车收车最早}：支线（往仁和方向）约 \textbf{22:00} 由新业路发出。
地铁晚得多，7 号线火车南站内环末班约 \textbf{23:27}，2 号线沿途各站末班普遍在 \textbf{23:00} 之后。\\
19:01 后若行程延误超过 1 小时，直接打车，不要赌末班车。\\
{\small\color{graytx} 各线时刻请以站台公告为准，本表按公开时刻表整理。}
\end{warnbox}

\begin{keybox}
\textbf{今日不安排签到与领物资}：抵达酒店已 20:35，CCF 签到点通常不会开放到深夜。
\textbf{CCFLinks 签到、领物资、系统调试统一排在次日赛前}（见 Day 2），今晚只做一件事，\textbf{吃好、早睡}。
\end{keybox}

% ---------- Day 2 ----------
\newpage
\daytitle{Day 2}{10 月 21 日（周三）· 比赛日}

\begin{center}
\begin{tabularx}{\textwidth}{@{}L{1.6cm}L{4.4cm}L{3.5cm}X@{}}
\toprule
\thd{时间} & \thd{事项} & \thd{交通方式} & \thd{地点 / 备注}\\
\midrule
07:00 & 起床、洗漱、早餐 & — & 闹钟设两个；酒店餐厅\textbf{吃饱}，12 小时比赛消耗大\\
07:35 & 清点随身物品 & — & 身份证、学生证、准考证、纸质资料、笔、水\\
07:45 & 大堂集合，出发前往赛场 & 步行 / 打车 & 约 1 公里，步行 15 分钟或打车起步价\\
08:05 & 抵达西华大学 & — & 按现场指示牌前往签到区\\
08:10 & \textbf{CCFLinks 扫码签到} & — & 出示准考证与身份证，核对参赛信息\\
08:20 & \textbf{领取参赛物资} & — & 胸牌、秩序册、纪念品等，\textbf{当场清点}再离开\\
08:35 & 系统调试 & — & 试编译、确认环境、摆放纸质资料\\
08:50 & 就座待考 & — & 检查键盘鼠标；入场物品检查（\textbf{无 U 盘}）\\
09:00 & \textbf{比赛开始} & — & 12 小时，\textbf{中途不离场}\\
12:00 & 午餐 & — & 往年发热盒饭，\textbf{仍建议自备干粮}\\
18:00 & 晚餐 & — & 往年发热盒饭，\textbf{仍建议自备干粮}\\
21:00 & \textbf{比赛结束} & — & 确认代码已提交，带走全部物品\\
21:15 & 返回酒店 & 步行 / 打车 & 约 1 公里\\
21:40 & 夜宵 & — & 红光镇周边\\
23:00 & 休息 & — & —\\
\bottomrule
\end{tabularx}
\end{center}

\begin{warnbox}
\textbf{签到与物资领取的确切时间官方尚未公布}，往年惯例为上午 8:00 前后开放。
上表按 08:05 到场预留，\textbf{出发前请留意 CCF 官方邮件与 CCSP 官网通知}，
若时间提前，把「大堂集合」整体前移即可。
\end{warnbox}

\begin{warnbox}
\textbf{补给是 12 小时比赛的关键}。\textbf{往年主办方午餐、晚餐都提供热盒饭}，
但\textbf{仍建议自备干粮}——排队领饭会打断做题节奏，手边有东西更稳妥。
建议随身携带：士力架 / 能量棒、咖啡、面包或饭团、不少于 1 升的饮水。
\end{warnbox}

\begin{okbox}
\textbf{赛前策略提醒}\\
$\bullet$ 5 题 500 分，一般难度递增。\\
$\bullet$ 纸质模板提前分册，考场上翻得快。\\
$\bullet$ 比赛 12 小时，建议每一段时间放松一次精神，避免后程思维僵化。\\
$\bullet$ \textbf{一定要学会骗分}，奖牌边缘争的往往不是完整过题，而是部分分：
把暴力、特殊性质、小数据分稳稳拿住。
\end{okbox}

% ---------- Day 3 ----------
\newpage
\daytitle{Day 3}{10 月 22 日（周四）· 颁奖典礼与 CNCC}

本日无比赛，\textbf{上午自由安排}，睡到自然醒也可以。只需记住一个节点：

\textbf{颁奖典礼 14:00--16:00，13:30 入场就座}，地点为\textbf{天府国际会议中心}
（成都市双流区蜀州路 3333 号，西博城地铁站 E 口步行 60 米）。

\begin{center}
\begin{tabularx}{\textwidth}{@{}L{2.3cm}L{1.9cm}L{2.8cm}X@{}}
\toprule
\thd{方案} & \thd{最晚出发} & \thd{耗时} & \thd{说明}\\
\midrule
打车 & \textbf{12:30} & 约 55 分钟 & 约 40 公里，120--140 元；\textbf{推荐}，最省心\\
地铁 6 号线 & \textbf{11:45} & 约 1 小时 40 分 & 一车直达西博城，无需二次换乘\\
\bottomrule
\end{tabularx}
\end{center}

\begin{okbox}
\textbf{CNCC2026 与颁奖典礼在同一片场馆}\\
\textbf{参赛选手免费参会}（10 月 21--24 日）。想听报告的话，
\textbf{上午先去 CNCC，下午原地参加颁奖}。
\end{okbox}

\begin{keybox}
\textbf{地铁 6 号线怎么走}\\
酒店 $\rightarrow$ 步行至\textbf{金土站} $\rightarrow$ 有轨电车蓉 2 号线支线（往新业路方向）
$\rightarrow$ \textbf{交大犀浦校区站}换乘 \textbf{地铁 6 号线}（往兰家沟方向）$\rightarrow$
\textbf{西博城站}下车 $\rightarrow$ \textbf{E 口}出站步行 60 米。
\end{keybox}

\begin{warnbox}
颁奖安排在天府新区，\textbf{与酒店分处城市两端}，单程 40 公里以上。
\end{warnbox}

% ---------- Day 3 ----------
\daytitle{Day 4}{10 月 23 日（周五）· 成都 $\rightarrow$ 合肥}\label{sec:day4}

上午同样自由安排，\textbf{只要在下面任一「最晚出发」之前出门}，就能赶上 11:51 的高铁。
两条路都按 \textbf{11:10 前抵达成都东站}倒推，留约 40 分钟安检进站。

\begin{center}
\begin{tabularx}{\textwidth}{@{}L{2.3cm}L{1.9cm}L{2.8cm}X@{}}
\toprule
\thd{方案} & \thd{最晚出发} & \thd{耗时} & \thd{说明}\\
\midrule
打车 & \textbf{10:20} & 约 50 分钟 & 约 30 公里，85--100 元；行李多或下雨直接打车\\
地铁 2 号线 & \textbf{09:35} & 约 1 小时 35 分 & 金土站 $\rightarrow$ 有轨电车 $\rightarrow$ 天河路换 2 号线（龙泉驿方向）$\rightarrow$ 成都东客站\\
\bottomrule
\end{tabularx}
\end{center}

\begin{keybox}
\textbf{出门前别忘了}\\
$\bullet$ 退房、结账，需要报销的记得索要发票；\\
$\bullet$ 重点检查充电器、充电宝、身份证等证件；\\
$\bullet$ 到站后按指示牌走，\textbf{检票口 A13}。
\end{keybox}

\begin{warnbox}
\textbf{车上历时 9 小时 50 分}：午餐、晚餐均在车上解决，\textbf{请提前备好食物与饮水}。
21:41 抵达合肥，当晚如需返校请提前安排交通。
\end{warnbox}

% ============================================================
\section{住宿信息}
% ============================================================

\begin{center}
\begin{tabularx}{\textwidth}{@{}L{3.0cm}X@{}}
\toprule
\thd{项目} & \thd{内容}\\
\midrule
酒店名称 & 繁洲花园大酒店（原丽天花园大酒店）\\
星级 & 四星级\\
地址 & 四川省成都市郫都区红光街道广场路北三段 88 号\\
电话 & 028-87922176 \quad/\quad 15308236858\\
距赛场 & 约 1 公里（部分来源标注 850 米，步行 12--18 分钟）\\
入住日期 & 10 月 20 日、21 日、22 日，共三晚\\
退房日期 & 10 月 23 日上午\\
\bottomrule
\end{tabularx}
\end{center}

\begin{keybox}
\textbf{入住当晚要问前台两件事}\\
1. 早餐几点开始、在哪一层\\
2. 10 月 23 日最晚退房时间，是否需要延迟退房
\end{keybox}

\begin{warnbox}
酒店电话存在多个号码，上表两个为较可靠的来源。\textbf{出发前建议再核实一次}，
并把号码存进手机通讯录，避免到达后临时查找。
\end{warnbox}

% ============================================================
\section{行前清单与提醒}\label{sec:remind}
% ============================================================

\subsection{随身证件与材料}

\begin{center}
\begin{tabularx}{\textwidth}{@{}L{0.9cm}L{4.2cm}X@{}}
\toprule
\thd{$\checkmark$} & \thd{物品} & \thd{说明}\\
\midrule
$\square$ & 身份证 & 乘车、入住、入场都要用，\textbf{时刻贴身保管}\\
$\square$ & 学生证 & 入场核验\\
$\square$ & 准考证 & \textbf{纸质版}，建议打印 2 份分开放\\
$\square$ & 纸质参考资料 & 今年\textbf{禁止电子设备}，模板必须打印装订\\
\bottomrule
\end{tabularx}
\end{center}

\begin{warnbox}
\textbf{出发前请在 12306 核对车票状态}，确认已出票、乘车人信息无误。
\end{warnbox}

\subsection{竞赛装备}

\begin{center}
\begin{tabularx}{\textwidth}{@{}L{0.9cm}L{4.2cm}X@{}}
\toprule
\thd{$\checkmark$} & \thd{物品} & \thd{说明}\\
\midrule
$\square$ & 键盘 / 鼠标 & 允许携带，无存储功能即可，须经检查\\
$\square$ & 有线耳机 & 允许携带\\
$\square$ & 笔、草稿本 & 竞赛用\\
$\square$ & 水杯 & 12 小时比赛，补水很重要\\
$\square$ & 干粮 & 士力架、能量棒、面包；考场是否供餐以现场为准\\
$\square$ & \textcolor{accent}{U 盘} & \textcolor{accent}{\textbf{不要带}}，今年明令禁止\\
\bottomrule
\end{tabularx}
\end{center}

\subsection{生活与气候}

\begin{keybox}
\textbf{成都 10 月下旬气候}\\
平均气温约 \textbf{13--19\,$^\circ$C{}}，白天 19\,$^\circ$C{} 左右、夜间 13\,$^\circ$C{} 左右，
\textbf{早晚温差明显}，多阴雨、湿度大。遇冷空气时最低可降到 10\,$^\circ$C{} 出头。\\[0.4em]
\textbf{穿衣建议}：长袖 T 恤或薄毛衣 + 夹克/风衣，外层可脱可穿；
带一把折叠伞；鞋子以防滑、耐走的运动鞋为宜。\\
\textbf{特别提醒}：要在考场里坐满 12 小时，\textbf{建议带一件薄外套}。
\end{keybox}

\begin{warnbox}
\textbf{饮食提醒}\\
成都菜\textbf{普遍偏辣}，点餐时\textbf{主动说明「外地人，少辣/不辣」}，店家一般会照顾口味。
比赛日前一晚\textbf{不要吃太辣、太油}，避免次日肠胃不适影响发挥。
\end{warnbox}

\subsection{其他}

\begin{itemize}[leftmargin=1.6em,itemsep=0.3em]
  \item 提前下载 \textbf{CCFLinks} App，并留意 CCF 官方发送的邮件。
  \item 出发前\textbf{告知亲友行程与目的地}，保持通讯畅通。
  \item 充电宝、数据线、转换头按需携带；酒店插座为标准国标。
  \item 常用药品（肠胃药、感冒药、晕车药）建议随身带一点。
  \item 竞赛期间\textbf{整体轻装上阵}，物品越少，检查越快、负担越轻。
\end{itemize}

\newpage

% ============================================================
\appendix
\section{资料来源}

本手册中的竞赛信息、交通路线与酒店信息来自以下来源，\textbf{出行前请以官方最新通知为准}：

\begin{itemize}[leftmargin=1.6em,itemsep=0.25em]
  \item CCF 官方通知：\href{https://www.ccf.org.cn/ccsp/tzgg/2026-09-09/931080.shtml}{2026 CCF CCSP 竞赛将于 10 月 21\textasciitilde 22 日举办}
  \item CCF 官方通知：\href{https://www.ccf.org.cn/ccsp/}{CCSP 竞赛官网}
  \item CCF 官方通知：\href{https://www.ccf.org.cn/ccsp/tzgg/2026-09-14/932132.shtml}{关于调整 2026 CCF CCSP 竞赛选手可携带资料规定的通知}（2026-09-14，CCSP 竞赛委员会）
  \item 成都地铁 2 号线、6 号线、7 号线线路资料
  \item 成都有轨电车蓉 2 号线支线线路资料
  \item 西华大学交通指南、酒店预订平台公开信息
\end{itemize}

\begin{warnbox}
本手册为\textbf{队伍内部出行参考}，其中交通耗时均为估算。
官方如有新的报到安排、场地变更或规则调整，\textbf{以官方通知为准}。
\end{warnbox}

\vfill
\begin{center}
{\color{mainlite}\rule{0.6\textwidth}{0.8pt}}\\[0.6em]
{\small\color{graytx} 祝旅途顺利，比赛取得理想成绩！}\\[0.4em]
{\color{mainlite}\rule{0.6\textwidth}{0.8pt}}
\end{center}

\end{document}
